\section{Exact isolation of the cross--row four--M\"obius form}
\label{sec:tpc41-physical-four}

The row--diagonal theorem leaves one terminal sector.  We now write it
without hiding its signs, masks, or alias pairs.

\subsection{Folded physical atoms}

For the minimal bank define
\begin{equation}
 u_{\alpha,j,0}=b_{\alpha,j,0},
 \qquad
 u_{\alpha,j,r}
 =b_{\alpha,j,r}+b_{\alpha,j,r-q}
 \quad(1\le r\le L).
 \label{eq:tpc41-physical-u}
\end{equation}
Hilbert--valued folded Parseval gives the exact left identity
\begin{align}
 \cE_L^{\rm fold}
 &:=\frac1q\sum_{\nu=0}^{q-1}
   \sum_{\gamma,j}
   \left|
    \sum_{\alpha,k}\omega^{\nu k}
    c^L_{\alpha,\gamma,j}b_{\alpha,j,k}
   \right|^2
 \notag\\
 &=\sum_{\gamma,j}\sum_{a=0}^{L}
   \left|
    \sum_\alpha c^L_{\alpha,\gamma,j}u_{\alpha,j,a}
   \right|^2.
 \label{eq:tpc41-physical-folded-energy}
\end{align}
Splitting the last square according to \(\alpha=\beta\) and
\(\alpha\ne\beta\), put
\begin{align}
 \cE_{{\rm same},L}
 &=\sum_{\gamma,j}\sum_{a=0}^{L}\sum_\alpha
   |c^L_{\alpha,\gamma,j}|^2|u_{\alpha,j,a}|^2,
 \label{eq:tpc41-physical-same-L}\\
 \cF_L
 &=\sum_{\gamma,j}\sum_{a=0}^{L}
   \sum_{\alpha\ne\beta}
   c^L_{\alpha,\gamma,j}
   \overline{c^L_{\beta,\gamma,j}}
   u_{\alpha,j,a}\overline{u_{\beta,j,a}}.
 \label{eq:tpc41-physical-FL}
\end{align}
The ordered row--pair sum is real, but we retain \(\operatorname{Re}\)
below to make clear that only an upper bound is used.  Define
\(\cE_{{\rm same},R}\) and \(\cF_R\) symmetrically.

\begin{theorem}[Exact terminal reduction]
\label{thm:tpc41-exact-terminal-reduction}
The minimal terminal bank satisfies
\begin{equation}
 \cE_L^{\rm fold}+\cE_R^{\rm fold}
 =\cE_{{\rm same},L}+\cE_{{\rm same},R}
  +\operatorname{Re}(\cF_L+\cF_R).
 \label{eq:tpc41-terminal-decomposition}
\end{equation}
The first two terms obey the endpoint target by
\eqref{eq:tpc41-physical-same-row-closed}.  Consequently the terminal
folded estimate
\begin{equation}
 \cE_L^{\rm fold}+\cE_R^{\rm fold}
 \ll X^{o(1)}Q^2JK_B
 \label{eq:tpc41-terminal-folded-target}
\end{equation}
is equivalent, up to the already proved diagonal term, to the one--sided
cross--row bound
\begin{equation}
 \boxed{
 \operatorname{Re}(\cF_L+\cF_R)
 \ll X^{o(1)}Q^2JK_B.}
 \label{eq:tpc41-terminal-cross-row-gate}
\end{equation}
\end{theorem}

\begin{proof}
Equation \eqref{eq:tpc41-terminal-decomposition} is the partition of the
ordered double row sum in \eqref{eq:tpc41-physical-folded-energy}.  The
same--row part is nonnegative and has the bound
\eqref{eq:tpc41-physical-same-row-closed}.  Hence
\eqref{eq:tpc41-terminal-cross-row-gate} implies
\eqref{eq:tpc41-terminal-folded-target}.  Conversely,
\(\operatorname{Re}(\cF_L+\cF_R)\) is the total energy minus a
nonnegative same--row energy, so \eqref{eq:tpc41-terminal-folded-target}
implies the same one--sided upper bound.
\end{proof}

This equivalence explains why estimating \(|\cF_L|+|\cF_R|\) is more than
is logically necessary.  A signed dispersion argument may exploit
cancellation between cross--row terms as long as it produces the upper
bound in \eqref{eq:tpc41-terminal-cross-row-gate}.

\subsection{The literal arithmetic expansion}

For \(a=0\), the only alias pair is \((k,k')=(0,0)\).  For
\(1\le a=r\le L\), expansion of
\(u_{\alpha,j,r}\overline{u_{\beta,j,r}}\) produces
\begin{equation}
 (k,k')\in
 \{(r,r),(r,r-q),(r-q,r),(r-q,r-q)\}.
 \label{eq:tpc41-physical-alias-pairs}
\end{equation}
The source factors inside \(c^L\) and the target atoms inside \(u\) show
that every summand with \(\alpha\ne\beta\) contains
\begin{equation}
 \boxed{
 \mu(d_\alpha)\mu(d_\beta)
 \mu(m_\alpha j+h+kM)
 \mu(m_\beta j+h+k'M).}
 \label{eq:tpc41-physical-four-mobius}
\end{equation}
It is multiplied by the two source logarithms, the two target logarithms,
the terminal cutoffs, the shared opposite--row masks, and the inherited
residue and smooth factors.  Formula
\eqref{eq:tpc41-physical-four-mobius} is therefore a description of the
actual coefficient kernel, not a replacement by an unweighted Chowla
model.

\subsection{Nonproportionality of every cross--row pair}

For fixed retained \(\alpha,\beta,k,k'\), define affine forms in \(j\) by
\begin{equation}
 L_{\alpha,k}(j)=m_\alpha j+h+kM,
 \qquad
 L_{\beta,k'}(j)=m_\beta j+h+k'M.
 \label{eq:tpc41-affine-forms}
\end{equation}

\begin{lemma}[Uniform algebraic nondegeneracy]
\label{lem:tpc41-affine-nonproportional}
If \(\alpha\ne\beta\), the two forms in
\eqref{eq:tpc41-affine-forms} are nonproportional for every retained
\(k,k'\).
\end{lemma}

\begin{proof}
Proportionality would force their coefficient determinant to vanish:
\begin{equation}
 m_\alpha(h+k'M)=m_\beta(h+kM).
 \label{eq:tpc41-affine-determinant-zero}
\end{equation}
Reducing modulo \(M\) and using \((h,M)=1\) gives
\(m_\alpha\equiv m_\beta\pmod M\).  The inherited row gap
\eqref{eq:tpc41-interface-row-gap} then gives
\(m_\alpha=m_\beta\).  Source separation forces
\(\alpha=\beta\), a contradiction.
\end{proof}

The lemma removes a local algebraic obstruction, but it does not import a
known affine Chowla theorem.  The slopes \(m_\alpha,m_\beta\) and
intercepts \(h+kM,h+k'M\) grow with \(X\), and the desired estimate sums
over row pairs with shared masks.  Even a logarithmic saving for every
fixed pair would not survive coefficient--blind absolute summation at the
required endpoint scale.

\subsection{Three distinct saving ledgers}

It is useful to distinguish three normalization problems:
\begin{center}
\begin{tabular}{@{}p{0.34\textwidth}ccc@{}}
\toprule
object & coherent scale & target scale & required gain\\
\midrule
raw equality terminal \(k=0\)
 & \(Q^3J\) & \(Q^3/J\) & \(J^2=Q/K_B\)\\
minimal folded or full trace terminal
 & \(Q^3JK_B\) & \(Q^2JK_B\) & \(Q\)\\
scalar length--\(K_B\) progression
 & \(MK_B^2\) & \(MK_B\) & \(K_B\)\\
\bottomrule
\end{tabular}
\end{center}
The same--row theorem removes the scalar long--gap expression from the
minimal folded physical gate.  It does not turn the cross--row row-energy
problem into the third line of the table.
